\subsection{Spectral measures}

In this section, we fix a complex Hilbert space E.

Recall that if A is a commutative unital stellar algebra, we denote by X(A) the compact topological space of its unital characters (cor. 1 of I, p. 29) and \(\mathcal{G}_{\mathrm{A}}\) the Gelfand transform of A (def. 5 of I, p. 7), which is an isometric isomorphism of unital involutive algebras from A onto \(\mathcal{C}(\mathrm{X}(\mathrm{A}))\) (th. 1 of I, p. 108). We will denote by \(\mathcal{H}_{\mathrm{A}}\) the inverse isomorphism.

Lemma 9. --- Let A be a commutative unital stellar subalgebra of \(\mathcal{L}(\mathrm{E})\) . For all \(x\) and \(y\) in E, the map \(\mu\) from \(\mathsf{X}(\mathrm{A})\) into \(\mathbf{C}\) defined by \(\mu(f) = \langle x | \mathcal{H}_{\mathrm{A}}(f)y \rangle\) for all \(f \in \mathcal{C}(\mathsf{X}(\mathrm{A}))\) is a bounded measure on the compact space \(\mathsf{X}(\mathrm{A})\) . It is positive if \(x = y\) . Its norm is \(\leqslant \|x\|\ \|y\|\) , with equality if \(x = y\) .

The map \(\mu\) is linear. For every function \(f\in\mathcal{C}(\mathsf{X}(\mathsf{A}))\) , we have

\[
| \mu (f) | \leqslant \| x \| \| y \| \| \mathcal {H} _ {\mathrm{A}} (f) \| = \| x \| \| y \| \| f \|,
\]

hence \(\mu\) is a bounded measure on \(\mathsf{X}(\mathsf{A})\) of norm \(\leqslant \| x\|\| y\|\) .

Suppose that \(x = y\) . For every positive function \(f \in \mathcal{C}(\mathsf{X}(\mathsf{A}))\) , the element \(\mathcal{H}_{\mathsf{A}}(f)\) is a positive element of A (example 1 of I, p. 115), hence a positive element of \(\mathcal{L}(\mathsf{E})\) , whence \(\mu(f) = \langle x | \mathcal{H}_{\mathsf{A}}(f)(x) \rangle \geqslant 0\) (prop. 8 of I, p. 138). This shows that \(\mu\) is a positive measure. Since \(\mu(1) = \|x\|^2\) the total mass of \(\mu\) is \(\|x\|^2\) (INT, III, p. 58, § 1, no. 8, cor. 2).

DEFINITION 2. --- Let A be a commutative unital stellar subalgebra of \(\mathcal{L}(\mathrm{E})\) . For all \(x\) and \(y\) in E, the linear form \(f \mapsto \langle x | \mathcal{H}_{\mathrm{A}}(f)y \rangle\) on \(\mathcal{C}(\mathrm{X}(\mathrm{A}))\) is called the spectral measure of \((x,y)\) relative to A. If \(x = y\) , we say that it is the spectral measure of \(x\) relative to A.

Remark. --- Let A be a commutative unital stellar subalgebra of \(\mathcal{L}(\mathrm{E})\) . For all \(x\) and \(y\) in E, denote \(\mu_{x,y}\) the spectral measure of \((x,y)\) relative to A. The map defined by \((x,y) \mapsto \mu_{x,y}\) from \(\mathrm{E} \times \mathrm{E}\) into \(\mathcal{M}^1(\mathrm{X}(\mathrm{A}))\) is sesquilinear.

Let \(u\) be a normal endomorphism of E and A the unital stellar subalgebra of \(\mathcal{L}(\mathrm{E})\) generated by \(u\) . It is commutative. The space \(\mathsf{X}(\mathrm{A})\) is identified with \(\mathrm{Sp}(u)\) by the mapping \(\chi \mapsto \chi(u)\) (lemma 10 of I, p. 109). The spectral measure \(\mu_{x,y}\) of \((x,y)\) relative to A is therefore identified with a measure on \(\mathrm{Sp}(u)\) , called the spectral measure of \((x,y)\) relative to \(u\) . For every function \(f \in \mathcal{C}(\mathrm{Sp}(u))\) , we then have

\[
\int_ {\mathrm{Sp} (u)} f d \mu_ {x, y} = \langle x | f (u) y \rangle
\]

according to def. 5 of I, p. 110.

